% sup3_03_S14a: Supplement part 3, step S14, first half: the vertex factors (Ba1 (3.14)-(3.15)) and the propagator expansion (Ba1 (3.44)-(3.45)) in the new background.

\subsection{S14, I: the expansion of the vertices and of the propagator in $\tilde B'$}\label{sup3:ssec:S14a}

\cite[p.~433]{Ba3}: ``We have $\tilde B=B_0+\tilde B'$, and we expand in $\tilde B'$ the expressions connected with each renormalized class of graphs using the formulas (I.3.14), (I.3.44), and (I.3.45). After this expansion we get a sum of expressions corresponding to some new big set of graphs, which can be decomposed into a sum of renormalized classes.'' The consumed form has three parts: (i) the expression of a class after the cube split equals a finite sum of expressions of graphs with scalar lines $G_k(\square,B_0)$, old vertices of the same types and localisations with $B_0$ in place of $\tilde B$ and extra external $\tilde B'$-legs, new vertices of the types \cite[(1.8)--(1.11)]{Ba3} and \cite[(1.13)--(1.15)]{Ba3} ``with external vector field $\tilde B'$ only'', at most $N_B$ of them on each old line, each with two scalar legs, no $A'$-leg and a factor $e_k$ per $\tilde B'$-leg, and remainder operators treated as external fields; (ii) the number of terms is bounded in terms of $\bar n$, $N_B$, $d$ only, and the new factors obey \cref{sup3:lem:S13}; (iii) the terms can be regrouped into renormalised classes. This subsection proves the two identities behind (i); \cref{sup3:ssec:S14b} proves (i)--(ii) and (iii).

\paragraph{The Taylor remainder.} For $z$ in the Lie algebra spanned by $q$ (real multiples of $q$, antisymmetric) put
\[
R_{\bar n+1}(z)\defeq(\bar n+1)\int_0^1(1-s)^{\bar n}e^{sz}\,ds,\qquad\text{so that}\qquad e^z=\sum_{j=0}^{\bar n}\frac{z^j}{j!}+\frac{z^{\bar n+1}}{(\bar n+1)!}R_{\bar n+1}(z)
\]
(Taylor's formula with integral remainder); $e^{sz}$ is orthogonal, so $\norm{R_{\bar n+1}(z)}\le(\bar n+1)\int_0^1(1-s)^{\bar n}ds=1$. Define, for a bond value $A$ and a contour value $A(\Gamma)$,
\[
F_{1,k}(A)\defeq\eta^{-1}\big(U(A)-1\big)=\sum_{j=1}^{\bar n}\eta^{j-1}\frac{(qe_kA)^j}{j!}+\eta^{\bar n}\frac{(qe_kA)^{\bar n+1}}{(\bar n+1)!}R_{\bar n+1}(\eta qe_kA),\qquad
U(A(\Gamma))-1=\sum_{j=1}^{\bar n}\frac{(qe_kA(\Gamma))^j}{j!}+\frac{(qe_kA(\Gamma))^{\bar n+1}}{(\bar n+1)!}R_{\bar n+1}(qe_kA(\Gamma)).
\]
These are \cite[(3.14)--(3.15)]{Ba1} (with the sign convention of this supplement; every identity below holds for either sign of the exponent). $F_{1,k}(A)$ is a sum of $\bar n+1$ terms, the $j$th homogeneous of degree $j$ in $e_kA$ with the factor $\eta^{j-1}$, and the last carrying the matrix $R_{\bar n+1}$ of norm $\le1$.

\begin{lemma}[the vertex factors]\label{sup3:lem:S141}
On $\square$, with $\tilde B=B_0+\tilde B'$ as in \cref{sup3:lem:S13}, for every bond $b\subset\square$ and contour $\Gamma\subset\square$:
\begin{enumerate}[label=\textup{(\alph*)},nosep,leftmargin=*]
\item $U(\tilde B_b)=U(B_{0,b})U(\tilde B'_b)$ and $U(\tilde B(\Gamma))=U(B_0(\Gamma))U(\tilde B'(\Gamma))$;
\item $U(\tilde B'_b)=1+\eta F_{1,k}(\tilde B'_b)$ and $U(\tilde B'(\Gamma))=1+[U(\tilde B'(\Gamma))-1]$, with the finite expansions displayed above;
\item $D^\eta_{\tilde B}\phi(b)=D^\eta_{B_0}\phi(b)+F_{1,k}(\tilde B'_b)U(B_{0,b})\phi(b_+)$;
\item[\textup{(c$'$)}] equivalently, in the form of \cite[(3.44)]{Ba1} with $(A,B)=(\tilde B',B_0)$: $D^\eta_{\tilde B}\phi(b)=U(\tilde B'_b)\big[D^\eta_{B_0}\phi(b)-F_{1,k}(-\tilde B'_b)\phi(b_-)\big]$;
\item $Q_k(\tilde B)\phi(y)=Q_k(B_0)\phi(y)+F_{2,k}(\tilde B',B_0)\phi(y)$, with $F_{2,k}(A,B)\phi(y)\defeq\sum_{x\in B^k(y)}\eta^d\,[U(A(\Gamma^{(k)}_{y,x}))-1]\,U(B(\Gamma^{(k)}_{y,x}))\phi(x)$.
\end{enumerate}
\end{lemma}

\begin{proof}
(a) The $U$ are exponentials of real multiples of the fixed matrix $q$; they commute and multiply, and $A\mapsto A(\Gamma)$ is additive. (b) is the definition of $F_{1,k}$ and the Taylor formula. (c) Insert (a)--(b) into $D^\eta_{\tilde B}\phi(b)=\eta^{-1}(U(\tilde B_b)\phi(b_+)-\phi(b_-))$: $U(\tilde B_b)\phi(b_+)=U(B_{0,b})\phi(b_+)+\eta F_{1,k}(\tilde B'_b)U(B_{0,b})\phi(b_+)$. (c$'$) Write $U(\tilde B_b)=U(\tilde B'_b)U(B_{0,b})$ and $\phi(b_-)=U(\tilde B'_b)U(-\tilde B'_b)\phi(b_-)$ with $U(-\tilde B'_b)=1+\eta F_{1,k}(-\tilde B'_b)$; then $\eta^{-1}(U(\tilde B_b)\phi(b_+)-\phi(b_-))=U(\tilde B'_b)\eta^{-1}\big(U(B_{0,b})\phi(b_+)-\phi(b_-)-\eta F_{1,k}(-\tilde B'_b)\phi(b_-)\big)$. (d) Insert (a) into the definition of $Q_k(\tilde B)$.
\end{proof}

\begin{lemma}[the propagator expansion and the vertex operators]\label{sup3:lem:S142}
Let $H_\square(A)=-\Delta^{\eta,N}_{A,\square}+m^2t^2+a_kP_k(A)$ be the Neumann operator on the sites of $\square$ and $G_k(\square,A)=H_\square(A)^{-1}$; both $H_\square(\tilde B)$ and $H_\square(B_0)$ are invertible since $m^2t^2>0$ (the identities below need nothing more; the uniform lower bound \cite[(1.8)]{Ba4} enters only the estimates of the fourth part). Put $G_0\defeq G_k(\square,B_0)$, $D_0\defeq D^\eta_{B_0}$, $Q_0\defeq Q_k(B_0)$, $(F\phi)(b)\defeq F_{1,k}(-\tilde B'_b)\phi(b_-)$, $F_2\defeq F_{2,k}(\tilde B',B_0)$, and $V_k\defeq H_\square(B_0)-H_\square(\tilde B)$. Then:
\begin{enumerate}[label=\textup{(\alph*)},nosep,leftmargin=*]
\item (\cite[(3.44)]{Ba1} with $(A,B)=(\tilde B',B_0)$.) $V_k=F^*D_0+D_0^*F-F^*F-a_k\big[F_2^*Q_0+Q_0^*F_2+F_2^*F_2\big]$, adjoints with respect to the $\eta^d$-weighted inner products on sites and bonds of $\square$.
\item (\cite[(3.45)]{Ba1}.) For every $N\ge0$, $G_k(\square,\tilde B)=\sum_{m=0}^{N}G_0(V_kG_0)^m+\cR_N$ with $\cR_N\defeq G_0(V_kG_0)^NV_kG_k(\square,\tilde B)$.
\item[\textup{(b$'$)}] The two-sided form, for $n_1,n_2\ge1$:
\[
G_k(\square,\tilde B)=\sum_{m<n_1}G_0(V_kG_0)^m+\sum_{m'<n_2}(G_0V_k)^{n_1}G_0(V_kG_0)^{m'}+(G_0V_k)^{n_1}G_k(\square,\tilde B)(V_kG_0)^{n_2}.
\]
Its last term is the analogue on $\square$ of the kernel \cite[(1.16)]{Ba3}; its main terms have at most $N_B\defeq n_1+n_2-1$ insertions of $V_k$.
\item Inserting the finite expansions of $F_{1,k}$ and $U(\tilde B'(\Gamma))-1$ into (a) writes $V_k$ as a sum of at most $6(\bar n+1)^2$ \emph{vertex operators}, each a bilinear form in $(\phi,\psi)$ with exactly two scalar legs and no $A'$-leg, every factor $\tilde B'$ accompanied by $e_k$, $B_0$ standing where $\tilde B$ stood in \cite[(1.8)--(1.15)]{Ba3} and $\tilde B'$ where $\tilde A$ stood:
\begin{itemize}[nosep,leftmargin=*]
\item the terms of $F^*D_0$, $D_0^*F$ are $\sum_b\eta^d\big[D_0\psi(b)\cdot\eta^{j-1}\tfrac{(-qe_k\tilde B'_b)^j}{j!}\phi(b_-)\big]$ and their adjoints, $1\le j\le\bar n$, with the $R$-term for $j=\bar n+1$: the forms \cite[(1.8)]{Ba3} at $n=0$, $n'=j$, and \cite[(1.9)]{Ba3};
\item the terms of $F^*F$ are $\sum_b\eta^d\eta^{j+j'-2}(e_k\tilde B'_b)^{j+j'}\big[\psi(b_-)\cdot c_{jj'}q^{j+j'}\phi(b_-)\big]$ with scalar coefficients $c_{jj'}$: the forms \cite[(1.10)]{Ba3}, and \cite[(1.11)]{Ba3} when an $R$-factor occurs. \emph{The terms with $j+j'$ odd cancel as operators}: the $(j,j')$ and $(j',j)$ terms are $(-1)^j\frac{(z)^{j+j'}}{j!j'!}$ and $(-1)^{j'}\frac{(z)^{j+j'}}{j!j'!}$ with $z=-\eta qe_k\tilde B'_b$, of opposite signs when $j+j'$ is odd. Equivalently $F_{1,k}(A)^TF_{1,k}(A)=\eta^{-2}(2-e^z-e^{-z})$ with $z=\eta qe_kA$, which is even in $z$. This is the parity condition ``$n+n'$ even'' of \cite[(1.10)]{Ba3}; the terms with an $R$-factor carry no parity, as \cite[(1.11)]{Ba3} carries none. A single odd $(j,j')$ term is a nonzero operator, so the cancellation is between the two orderings, not termwise;
\item the terms with $F_2$ are block sums $\sum_y$ of products of a block factor \cite[(1.13)]{Ba3} at $B_0$ with a factor \cite[(1.14)]{Ba3} or \cite[(1.15)]{Ba3} at $n=0$, and of two such factors: the products \cite[(1.18)]{Ba3}.
\end{itemize}
The $\eta$-powers per leg are those of \cite[(1.8)--(1.15)]{Ba3}: $\eta^{j-1}$ at a term of type \cite[(1.8)]{Ba3} and $\eta^{j+j'-2}$ at a term of type \cite[(1.10)]{Ba3}.
\end{enumerate}
\end{lemma}

\begin{proof}
(a) By \cref{sup3:lem:S141}(c$'$) and (d), as operators on the fields of $\square$, $D^\eta_{\tilde B}=U(\tilde B')(D_0-F)$ with $U(\tilde B')$ the bondwise orthogonal multiplication by $U(\tilde B'_b)$, and $Q_k(\tilde B)=Q_0+F_2$. Hence $(D^\eta_{\tilde B})^*D^\eta_{\tilde B}=(D_0-F)^*(D_0-F)=D_0^*D_0-D_0^*F-F^*D_0+F^*F$ and $P_k(\tilde B)=(Q_0+F_2)^*(Q_0+F_2)$. Subtracting $H_\square(\tilde B)$ from $H_\square(B_0)$, the mass terms cancel and $V_k$ is the stated bracket. (The factorisation of \cref{sup3:lem:S141}(c), $D^\eta_{\tilde B}=D_0+\mathfrak F$ with $(\mathfrak F\phi)(b)\defeq F_{1,k}(\tilde B'_b)U(B_{0,b})\phi(b_+)$, gives the different exact form $V_k=-[D_0^*\mathfrak F+\mathfrak F^*D_0+\mathfrak F^*\mathfrak F]-a_k[F_2^*Q_0+Q_0^*F_2+F_2^*F_2]$; the two brackets are equal as operators but not term by term. The form (a) is the one whose terms are the printed vertices.)

(b) $G_0^{-1}-G_k(\square,\tilde B)^{-1}=V_k$ gives the second resolvent identity $G_k(\square,\tilde B)=G_0+G_0V_kG_k(\square,\tilde B)$; iterate it $N$ times.

(b$'$) Apply the identity on the other side, $G_k(\square,\tilde B)=G_0+G_k(\square,\tilde B)V_kG_0$, $n_2$ times to the factor $G_k(\square,\tilde B)$ of the remainder of (b) at $N=n_1-1$, i.e.\ to $(G_0V_k)^{n_1}G_k(\square,\tilde B)$.

(c) Distribute (a) over the $\bar n+1$ terms of $F$ and of $F_2$: at most $2(\bar n+1)+(\bar n+1)^2$ terms from the three $D$-terms and the same from the three $Q$-terms, at most $6(\bar n+1)^2$ in all. Each term has one $\phi$ and one $\psi$ leg (two scalar legs) and no $A'$. The forms are read off: $(F\phi)(b)=F_{1,k}(-\tilde B'_b)\phi(b_-)$ has the undifferentiated leg at $b_-$ and no $U(B_0)$ factor, so $F^*D_0$ pairs $D_0\psi(b)$ with $q^j\phi(b_-)$, the printed shape ``$[(D^\eta_{\tilde B}\phi')(b)\cdot q^{n+n'}\phi'(b_-)]$'' of \cite[p.~413]{Ba3}; $F^*F$ pairs $\psi(b_-)$ with $\phi(b_-)$ through $F_{1,k}(-\tilde B'_b)^TF_{1,k}(-\tilde B'_b)$, whose $(j,j')$ term is $\eta^{j+j'-2}\frac{((-z)^T)^j(-z)^{j'}}{j!j'!}$ with $z=\eta qe_k\tilde B'_b$, and $(q^T)^j=(-1)^jq^j$ gives the parity statement. The $F_2$ terms are the definition of $F_{2,k}$ expanded.
\end{proof}

\begin{remark}\label{sup3:rem:S14a}
The identities of this subsection are exact finite-dimensional identities, valid for any invertible $H_\square(\tilde B)$, $H_\square(B_0)$ and any bond field $\tilde B'$; the regularity of \cref{sup3:lem:S13} is used only in \cref{sup3:lem:S143}(v) and in the estimates of the fourth part. \cite[(3.47)--(3.53)]{K1} prints the same expansion on the torus.
\end{remark}
